\documentclass[12pt,a4paper]{book}

\usepackage[T1]{fontenc}
\usepackage{amsmath,amsthm,amssymb,mathtools}
\usepackage{geometry}
\usepackage{microtype}
\usepackage{setspace}
\usepackage{titlesec}
\usepackage{titletoc}
\usepackage{hyperref}
\usepackage{enumitem}
\usepackage{booktabs}
\usepackage{fancyhdr}
\usepackage{emptypage}

\geometry{
  a4paper,
  top=2.8cm, bottom=2.8cm,
  inner=3.2cm, outer=2.4cm,
  headheight=14pt
}

\onehalfspacing

\hypersetup{
  colorlinks=true,
  linkcolor=black,
  citecolor=black,
  urlcolor=black,
  pdftitle={Foundations of Distinction Geometry},
  pdfauthor={Flyxion}
}

% Theorem environments
\newtheorem{theorem}{Theorem}[chapter]
\newtheorem{lemma}[theorem]{Lemma}
\newtheorem{corollary}[theorem]{Corollary}
\newtheorem{proposition}[theorem]{Proposition}
\newtheorem{definition}[theorem]{Definition}
\theoremstyle{remark}
\newtheorem{remark}{Remark}[chapter]
\newtheorem{example}{Example}[chapter]
\theoremstyle{definition}
\newtheorem{principle}{Principle}[chapter]
\newtheorem{conjecture}{Conjecture}[chapter]

% Chapter title formatting
\titleformat{\chapter}[display]
  {\normalfont\Large\bfseries}
  {\chaptertitlename\ \thechapter}
  {16pt}
  {\huge}

\titleformat{\section}{\large\bfseries}{\thesection.}{0.6em}{}
\titleformat{\subsection}{\normalsize\bfseries}{\thesubsection.}{0.5em}{}

% Headers
\pagestyle{fancy}
\fancyhf{}
\fancyhead[LE]{\small\itshape\leftmark}
\fancyhead[RO]{\small\itshape\rightmark}
\fancyfoot[C]{\thepage}
\renewcommand{\headrulewidth}{0.4pt}

%----------------------------------------------------------------------
\begin{document}
%----------------------------------------------------------------------

\frontmatter

%--- Title page
\begin{titlepage}
\centering
\vspace*{3cm}
{\Huge\bfseries Foundations of Distinction Geometry}\\[1em]
{\large\itshape A Unified Catalogue of the Flyxion Research Program}\\[3em]
{\Large Flyxion}\\[0.8em]
{\normalsize Independent Researcher}\\[3em]
{\normalsize\today}\\
\vfill
{\small
\textit{A unified framework showing how distinctions, admissibility,\\
repair, projection, and coordination arise from a single architecture.}}
\end{titlepage}

%--- Epigraph page
\cleardoublepage
\vspace*{6cm}
\begin{center}
\itshape
``What exists?''\\[0.8em]
$\longrightarrow$\\[0.8em]
``What distinctions survive projection?''
\end{center}
\vspace*{3cm}
\begin{center}
\itshape
``A system can only act, remember, predict, coordinate, or persist\\
to the extent that it maintains distinctions.''
\end{center}

%--- Table of contents
\cleardoublepage
\tableofcontents

%--- Preface
\cleardoublepage
\chapter*{Preface}
\addcontentsline{toc}{chapter}{Preface}

This document is both a catalogue and an argument. The catalogue records
the frameworks, theories, formal systems, and computational artefacts
developed under the name Flyxion over roughly a decade of independent
research. The argument is that they are not separate inventions. They are
cross-sections through a single architecture, and that architecture has
a name: distinction geometry.

The central observation is simple. A system can only act, remember,
predict, coordinate, or persist to the extent that it maintains
distinctions. When distinctions collapse --- through noise, compression,
forgetting, institutional drift, physical coarse-graining, or the
inherent many-to-one character of any observation --- something is lost
that cannot be recovered without additional information. The geometry of
that loss, and the geometry of what survives, is the subject of this
program.

The frameworks documented here include: \textbf{RSVP} (Relativistic
Scalar-Vector Plenum, a field-theoretic cosmology grounded in
constraint-first dynamics); \textbf{CLIO} (constraint-layered
inference and projection theory); \textbf{TARTAN} (trajectory-aware
recursive tiling with annotated noise); \textbf{HYDRA} (coordination
geometry and collective admissibility as categorical colimit);
\textbf{Spherepop} (an irreversible event calculus with Pop, Refuse,
Bind, and Collapse operators); the \textbf{Admissibility Program}
(distinguishability geometry and reachability-based ontology); the
\textbf{Repair Theory} (repair as a primitive prior to optimization);
\textbf{Distinction Geometry} (Jensen--Shannon interface metrics and
the curvature of distinguishability); the \textbf{Preference Field
Program}; and the \textbf{Ecology of Distinctions} (a 31-chapter
academic monograph on distinction theory, reachability, and
admissibility geometry).

Earlier works often asked: \textit{What structures exist in the internal
space?} The current program asks: \textit{What can be known about the
internal space from observations?} And then takes the further step:
\textit{Perhaps the geometry of that limitation is itself the primary
object.}

This document is organized so that a reader beginning with Chapter~1
arrives at quadratic gravity, coordination geometry, the emergence of
quantum probability, and the geometry of scientific revolutions as
consequences of the same underlying architecture, rather than as dozens
of independent projects.

\bigskip\noindent
\textit{Flyxion}\\
Independent Researcher

\mainmatter

%======================================================================
\part{The Primitive Problem}
%======================================================================

\chapter{Why Distinctions Matter}

\section{The Central Observation}

A system can only act, remember, predict, coordinate, or persist to the
extent that it maintains distinctions. This is not a claim about
information theory, or about physics, or about cognition in particular.
It is a claim about what any system must do to be a system at all.

Consider action. To act differently in different circumstances, a system
must be able to distinguish those circumstances. A thermostat that cannot
distinguish hot from cold cannot regulate temperature. A market that cannot
distinguish solvent from insolvent firms cannot allocate capital. A nervous
system that cannot distinguish threat from safety cannot protect the
organism it inhabits. Action requires discrimination, and discrimination
requires maintained distinction.

Consider memory. To remember is to preserve a distinction between what was
and what is, between an earlier state and the current one. A system with
no distinctions has no memory, not because it lacks storage, but because
it cannot mark one time as different from another. The content of memory
is the set of distinctions that have been preserved across time.

Consider prediction. To predict is to distinguish possible futures from one
another and from the present. A prediction that cannot distinguish outcomes
carries no information about what will happen. The quality of a predictive
model is precisely the quality of the distinctions it can maintain across
time and circumstance.

Consider coordination. For two systems to coordinate, they must share
distinctions: they must agree on what counts as the same and what counts
as different. Language is a technology for sharing distinctions. Law is
a technology for enforcing them. Mathematics is a technology for making
them precise.

Consider persistence. A system persists when its identity-constituting
distinctions are maintained against the pressures that would collapse them.
An organism persists when it maintains the distinction between self and
environment. An institution persists when it maintains the distinctions
between its roles, its procedures, and its goals. A physical structure
persists when the forces that would homogenize it are outweighed by the
forces that maintain its differentiation.

\begin{principle}[Distinction Priority]
Distinctions are prior to information, entropy, memory, computation,
agency, knowledge, and coordination. Each of these concepts can be
derived from the notion of maintained distinction; none of them can be
defined without it.
\end{principle}

\section{What Distinction Collapse Looks Like}

The collapse of distinctions is one of the most pervasive phenomena in
nature and in human affairs, and it takes many forms.

\textit{Physical collapse.} Thermodynamic equilibrium is the state in
which all thermodynamic distinctions within a system have been eliminated:
temperature, pressure, and chemical potential are uniform everywhere.
The approach to equilibrium is the approach to distinction collapse.
Entropy, in this reading, is a measure of how many distinctions have
been lost: how many microstates are compatible with the current macrostate,
how large the preimage of the current observation.

\textit{Cognitive collapse.} Forgetting is the collapse of temporal
distinctions: the inability to mark what was as different from what is.
Cognitive biases are often distinction collapses: the availability
heuristic collapses the distinction between frequency and memorability;
the halo effect collapses the distinction between one trait and all
others; confirmation bias collapses the distinction between evidence and
expectation.

\textit{Institutional collapse.} Goodhart's Law is a distinction collapse:
when a measure becomes a target, the distinction between the measure and
the thing it was measuring collapses. Goal drift is a distinction collapse:
the distinction between the stated purpose of an institution and its
actual behavior erodes over time. Bureaucratic ossification is a
distinction collapse: the distinction between the form and the substance
of a procedure disappears, leaving only the form.

\textit{Scientific collapse.} Paradigm exhaustion, in Kuhn's sense, is
a distinction collapse: the reigning framework can no longer distinguish
anomalies from noise, can no longer separate what it predicts from what
it observes. A scientific revolution is a distinction reorganization:
new categories are introduced that make previously invisible differences
visible.

\textit{Computational collapse.} Information loss in computation is a
distinction collapse. Irreversible operations --- AND, OR, NAND --- map
multiple input states to a single output state, collapsing the distinction
between inputs. The thermodynamic cost of irreversible computation is
the energy required to erase the record of that collapsed distinction.

\section{Why This is the Right Starting Point}

The history of the foundational sciences is a history of discovering that
something thought to be primitive is in fact derived from something more
basic. Newtonian mechanics was derived from the calculus of variations.
Thermodynamics was derived from statistical mechanics. Information theory
was derived from probability. Probability is currently being rederived
from admissibility and projection.

The Flyxion program proposes that distinctions are the next primitive:
that information, entropy, probability, and admissibility are all derived
from the more basic notion of maintained distinction. This is not a
speculative claim. The derivations are largely in hand; they are
documented in the subsequent chapters of this work.

What is primitive is not an object, a state, a field, or a wave function.
What is primitive is the capacity to tell things apart.

\chapter{Distinction Spaces}

\section{Formal Definition}

\begin{definition}[Distinction Space]
A \emph{distinction space} is a pair $(X, \sim)$ where $X$ is a set
and $\sim$ is an equivalence relation on $X$. An element $x \in X$ is
a \emph{history}; the equivalence class $[x]_\sim$ is the \emph{observable
class} of $x$. Two histories $x, x' \in X$ are \emph{indistinguishable}
if $x \sim x'$ and \emph{distinguishable} otherwise.
\end{definition}

Every equivalence relation on $X$ induces a projection
\begin{equation}
  \Pi : X \to X/{\sim}
\end{equation}
onto the quotient set of equivalence classes. The projection $\Pi$ records
everything that can be observed; the fiber $\Pi^{-1}([x])$ records everything
that cannot.

\begin{remark}
This is a very general setup. When $X$ is the microstate space of a
thermodynamic system and $\sim$ is thermodynamic coarse-graining, we
recover the foundations of statistical mechanics. When $X$ is the space
of neural activation patterns and $\sim$ is behavioral equivalence, we
recover a framework for cognitive science. When $X$ is the space of gauge
field configurations and $\sim$ is gauge equivalence, we recover the fiber
bundle structure of Yang--Mills theory. The apparatus of distinction spaces
is not a metaphor across these domains; it is the same mathematical object.
\end{remark}

\section{Distinguishability Classes and Their Properties}

The quotient $X/{\sim}$ is the space of \emph{distinguishability classes}:
the collection of all things that can be told apart. It inherits whatever
structure $X$ carries that is compatible with $\sim$. When $X$ is a
topological space and $\sim$ is an open equivalence relation, $X/{\sim}$
is a topological space. When $X$ is a measure space and $\sim$ is
measurable, $X/{\sim}$ is a measure space.

\begin{definition}[Refinement and Coarsening]
A distinction space $(X, \sim')$ is a \emph{refinement} of $(X, \sim)$
if $x \sim' y \Rightarrow x \sim y$: finer equivalences are available.
It is a \emph{coarsening} if $x \sim y \Rightarrow x \sim' y$: fewer
distinctions are maintained. Coarsening is projection; refinement is
distinction creation.
\end{definition}

\begin{definition}[Representational Entropy]
Given a distinction space $(X, \sim, \mu)$ with reference measure $\mu$,
the \emph{representational entropy} of the equivalence relation $\sim$
is
\begin{equation}
  H(\sim \mid \mu)
  = -\sum_{[x] \in X/\sim} \mu([x]) \log \mu([x]).
\end{equation}
This measures the uncertainty about which equivalence class a randomly
drawn history belongs to. It is maximised when all classes are equally
likely (maximum distinction preservation) and minimised when all histories
are in a single class (complete distinction collapse).
\end{definition}

\section{Projection, Collapse, and Compression}

The three fundamental operations on distinction spaces are projection,
collapse, and compression. They are related but distinct.

\textit{Projection} is the map $\Pi : X \to X/{\sim}$. It sends each
history to its equivalence class. It is the operation of observation:
it records what is visible and discards what lies in the fiber.

\textit{Collapse} is a change in the equivalence relation that coarsens
it: $\sim$ is replaced by a coarser $\sim'$. It corresponds to the
process by which a system loses the ability to distinguish things it
previously could distinguish. Forgetting is cognitive collapse.
Equilibration is thermodynamic collapse. Goal drift is institutional
collapse.

\textit{Compression} is a deliberate coarsening that preserves a
specified subset of distinctions while collapsing others. Efficient
compression preserves the distinctions that matter for a given purpose
while eliminating those that do not. The theory of compression is
therefore the theory of which distinctions matter for which purposes ---
a question that connects compression theory to decision theory, value
theory, and the theory of admissibility.

\chapter{The Lifecycle of Distinctions}

\section{Creation}

Distinctions are created when a system acquires the capacity to tell
things apart that it previously could not. This happens through:

\textit{Differentiation.} A homogeneous system spontaneously breaks into
distinguishable regions. Phase transitions in physics, cell differentiation
in biology, and the emergence of social roles in institutions are all
examples of distinction creation through differentiation.

\textit{Learning.} A cognitive system updates its projection $\Pi$ to
distinguish stimuli that it previously conflated. Perceptual learning ---
the ability to distinguish wines, musical intervals, or faces with
increasing precision --- is distinction creation.

\textit{Measurement.} The construction of a measuring instrument creates
distinctions by coupling a physical system to a record-keeping apparatus
whose states correspond to the distinctions being introduced. The history
of science is substantially the history of constructing instruments that
create previously unavailable distinctions.

\textit{Language.} The introduction of a new word or concept creates a
distinction by providing a label that separates what was previously
unseparated. Technical vocabulary is distinction-creation technology.

\section{Preservation}

A distinction is preserved when the processes that would collapse it ---
noise, entropy, drift, forgetting --- are outweighed by the processes
that maintain it. Preservation requires active work.

The energy cost of maintaining a distinction is a central concern of
Landauer's principle: erasing one bit of information requires a minimum
of $k_BT\ln 2$ joules. Maintaining a distinction against thermal noise
requires a comparable investment. This connects distinction preservation
to thermodynamics and to the energetics of life.

The social cost of maintaining a distinction is equally real. Legal
distinctions require enforcement. Linguistic distinctions require use.
Institutional distinctions require active reproduction through practice.
A distinction that is not exercised tends to collapse.

\section{Repair}

Repair is the restoration of a distinction that has been partially or
fully collapsed. It is the process of recovering information about the
generating history from its observable image, when that image is
insufficient to determine the history uniquely.

\begin{principle}[Repair Priority]
Repair is prior to optimization, prediction, and intelligence. A system
that cannot repair distinctions that have been collapsed cannot learn from
experience, cannot correct errors, and cannot adapt to novelty. Repair
is the capacity that makes all other cognitive capacities possible.
\end{principle}

Repair takes many forms: error correction in coding theory, memory
retrieval in cognitive science, paradigm revision in epistemology,
immune response in biology, constitutional amendment in political theory.
All of these are instances of the same operation: the partial restoration
of a collapsed distinction from whatever information survives.

Repair Theory is developed in Part~IV of this work.

\section{Loss}

Distinction loss is irreversible in the precise sense established by
Distinction Geometry: when the projection $\Pi$ is non-injective and the
lost information is not recorded elsewhere, no reconstruction map can
perfectly recover the original history. The minimum reconstruction defect
is positive.

This irreversibility is not a failure of ingenuity. It is a theorem. The
theorem is what gives the second law of thermodynamics its force: the
coarse-graining projection from microstates to macrostates is non-injective,
and the minimum reconstruction defect grows with time under any
measure-preserving dynamics.

\section{Reorganization}

Reorganization is the replacement of one distinction space $(X, \sim)$
by another $(X, \sim')$, where $\sim'$ is neither a refinement nor a
coarsening of $\sim$: it makes distinctions that $\sim$ did not make and
collapses distinctions that $\sim$ preserved.

Scientific revolutions are distinction reorganizations. The Copernican
revolution replaced a distinction space in which the earth and the heavens
were distinguished by a distinction space in which the earth is a planet
like any other. The quantum revolution replaced a distinction space in
which position and momentum were simultaneously sharp by one in which
they are fundamentally conjugate.

Conceptual change, paradigm shifts, learning, and creativity are all,
in the present framework, instances of distinction reorganization. They
are the topic of Chapter~21.

%======================================================================
\part{Information and Projection}
%======================================================================

\chapter{Information as Distinction}

\section{Deriving Information Theory}

The standard foundations of information theory take probability as
primitive and derive entropy and information from it. The distinction
geometry program reverses this: distinctions are primitive, and
information is derived from them.

\begin{definition}[Information as Maintained Separation]
The \emph{information content} of a distinction between $x$ and $y$
in a distinction space $(X, \sim, \mu)$ is the negative log-probability
of the event that a randomly drawn pair $(x',y')$ belongs to the same
equivalence class:
\begin{equation}
  I(x \not\sim y)
  = -\log P\bigl([x]_\sim \neq [y]_\sim\bigr).
\end{equation}
The \emph{entropy} of the distinction space is the expected information
content: $H(\sim \mid \mu) = \mathbb{E}[I]$.
\end{definition}

Under this derivation, Shannon entropy emerges as the expected
distinguishability of a randomly drawn pair: the amount of distinction
that the equivalence relation $\sim$ maintains on average. High entropy
means many distinctions are maintained; low entropy means few.

\begin{theorem}
Shannon entropy $H(P) = -\sum_i p_i\log p_i$ equals the expected
information content of a uniform-weight distinction space $(X, \sim_P)$
where $\sim_P$ is the equivalence relation induced by the probability
distribution $P$.
\end{theorem}

\begin{proof}[Proof sketch]
The probability that a random pair $(x, y)$ drawn i.i.d.\ from $P$
falls in the same class is $\sum_i p_i^2$. The expected information
content of the distinction between them is $-\log(1 - \sum_i p_i^2)$
for distinguishable pairs. Under the standard information-theoretic
derivation from the axioms of entropy, the result is $H(P)$ as the
unique measure satisfying continuity, symmetry, and the chain rule for
conditional distinctions.
\end{proof}

\section{Compression as Distinction Reduction}

Lossless compression is the identification of a coarser representation
that preserves all distinctions relevant to a specified purpose. Lossy
compression deliberately collapses some distinctions to achieve a more
compact representation. The rate--distortion tradeoff is the tradeoff
between the number of distinctions preserved and the cost of preserving them.

The mutual information $I(X; Y)$ between two random variables is the
average amount of distinction about $X$ maintained by $Y$: the reduction
in uncertainty about $X$ given $Y$. In distinction geometry terms, it is
the expected decrease in fiber size induced by conditioning on $Y$.

\section{The Data Processing Inequality as a Distinction Theorem}

\begin{theorem}[Data Processing as Distinction Monotonicity]
If $X \to Y \to Z$ is a Markov chain (each subsequent variable is
a function of the previous one), then $I(X;Z) \leq I(X;Y)$. Processing
can only reduce the distinctions maintained about $X$.
\end{theorem}

The data processing inequality is, in distinction geometry terms, the
statement that projection is monotone: every additional projection
between the history and the observation can only collapse more
distinctions. It is impossible to create distinctions about $X$ by
processing $Y$.

\chapter{Projection Geometry: The CLIO Framework}

\section{Overview}

The CLIO framework (Constraint-Layered Inference and Projection) provides
the formal architecture for projection geometry. It was developed to
address a gap in the existing frameworks for inference and learning: the
failure to distinguish between what a system can infer from observations
(the observable manifold $M_{\mathrm{obs}}$) and what structures the
system projects onto the world (the internal manifold $M_{\mathrm{int}}$).

CLIO identifies three layers in any cognitive or inferential system:
\begin{enumerate}[label=(\roman*)]
  \item The \emph{constraint layer}: the set of admissibility conditions
        that a representation must satisfy to count as legitimate.
  \item The \emph{projection layer}: the map $\pi : M_{\mathrm{int}} \to
        M_{\mathrm{obs}}$ that connects internal representations to
        observable quantities.
  \item The \emph{inference layer}: the process by which the system
        updates its internal representation given new observations, i.e.,
        attempts to invert $\pi$.
\end{enumerate}

\section{The CLIO Projection}

The central object of CLIO is the projection
\begin{equation}
  \pi : M_{\mathrm{int}} \to M_{\mathrm{obs}},
\end{equation}
which maps each internal world-model to the set of observations it
predicts. Two world-models are \emph{projection-equivalent} if they
predict the same observations: $m \sim m' \iff \pi(m) = \pi(m')$.

\begin{principle}[CLIO Projection Sufficiency]
No condition on the internal model $m \in M_{\mathrm{int}}$ can be
physically or epistemically necessary unless it changes the projection
$\pi(m)$. Constraints belong on $M_{\mathrm{obs}}$, not on
$M_{\mathrm{int}}$.
\end{principle}

This principle is the epistemological foundation for the critique of
latent fundamentalism (Chapter~26) and the Observational--Interventional
Separation theorem.

\section{Projection Failure and World-Model Mismatch}

A \emph{projection failure} occurs when the internal model $m$ and the
observation $o$ are inconsistent: $o \notin \pi(m)$ or $\pi(m) \notin
\mathcal{A}_O$. Projection failures come in two kinds:

\textit{Structural failure}: the internal model's projection misses the
actual observation entirely. The model cannot explain what happened.

\textit{Admissibility failure}: the internal model's projection is
consistent with the observation but violates an admissibility constraint.
The model explains what happened but in a way that is not admissible
(e.g., violates conservation laws, causality, or computational constraints).

Repair Theory (Part~IV) is concerned primarily with structural failure:
how to update a world-model whose projection no longer includes the
current observation.

\section{Abstraction as Coarsening}

Abstraction, in CLIO, is the operation of coarsening the internal model:
replacing a fine-grained $M_{\mathrm{int}}$ with a coarser
$M_{\mathrm{int}}'$ in which formerly distinct models are identified.
Good abstraction preserves the distinctions that matter for prediction
while collapsing those that do not.

The theory of good abstraction is therefore the theory of which
distinctions in $M_{\mathrm{int}}$ are tracked by distinctions in
$M_{\mathrm{obs}}$: which internal distinctions have observable
consequences. This is exactly the question answered by the CLIO projection
$\pi$: the distinctions that matter are those not collapsed by $\pi$.

\section{Latent Fundamentalism and Its Critique}

\emph{Latent fundamentalism} is the error of granting ontological status
to internal coordinates merely because they appear in a successful model.
It arises when the projection $\pi$ is forgotten and the internal model
is treated as if it were the observable world.

The error is pervasive:
\begin{itemize}
  \item In quantum mechanics: treating the wave function as a physical
        object rather than a tool for generating predictions about
        observations.
  \item In neuroscience: treating neural activation patterns as the
        content of experience rather than as projections of that content
        onto a measurable substrate.
  \item In economics: treating utility functions as real preferences
        rather than as models of revealed preferences under projection.
  \item In institutional theory: treating the organizational chart as
        the institution rather than as a projection of actual power and
        decision-making structures.
\end{itemize}

In each case, the correction is the same: identify the projection $\pi$,
restrict admissibility conditions to $M_{\mathrm{obs}}$, and treat
$M_{\mathrm{int}}$ as a representational instrument rather than an
ontological object.

\chapter{Interface Theory}

\section{The Observable--Internal Distinction}

Interface theory is the formal development of the distinction between
what is observable and what is internal. Its central claim is that the
observable interface of a system --- the distribution $P = \pi_*\mu$ on
$M_{\mathrm{obs}}$ --- is the only thing about the system that is
empirically accessible.

\begin{definition}[Observable Interface]
Given a projection setting $(M, O, \Pi, \mu)$, the \emph{observable
interface} of the system is the pushforward distribution
$P = \Pi_*\mu \in \mathcal{M}$, where $\mathcal{M}$ is the space of
probability distributions on $O$.
\end{definition}

Two systems with the same observable interface are \emph{interface-equivalent}:
they cannot be distinguished by any observation, however fine-grained.

\section{Observational--Interventional Separation}

The Observational--Interventional Separation theorem establishes that the
observable interface underdetermines the internal structure:

\begin{theorem}[Observational--Interventional Separation]
The map from internal structures to observable interfaces is not injective:
multiple distinct internal structures $(M^{(1)}, \Pi_1, \mu_1)$ and
$(M^{(2)}, \Pi_2, \mu_2)$ may share the same observable interface
$P^{(1)} = P^{(2)}$.
\end{theorem}

This theorem has immediate practical consequences. It establishes that
no experiment --- however carefully designed --- can uniquely determine
the internal structure of a system from its observable interface. The
choice of internal model is therefore always underdetermined by evidence.

The principle of ontological restraint follows: admissibility conditions
should be imposed on observable interfaces, not on internal structures.
Internal structures are representational tools; they should be evaluated
by their predictive accuracy and computational utility, not by their
ontological plausibility.

\section{Projection Sufficiency Principle}

\begin{principle}[Projection Sufficiency]
Physical, cognitive, and institutional predictions depend only on the
observable interface $P = \Pi_*\mu$. No condition on the internal
structure $(M, \Pi, \mu)$ is necessary unless it changes the observable
interface.
\end{principle}

This principle unifies the ghost problem in quantum gravity, the
underdetermination problem in philosophy of science, the model selection
problem in statistics, and the legibility problem in institutional theory.
In each case, the operative principle is the same: constraints belong on
the interface, not on the internal structure.

\section{Interface-Native Metrics: The Jensen--Shannon Distance}

The appropriate metric for comparing observable interfaces is the
Jensen--Shannon distance, developed in detail in Distinction Geometry
(Chapter~27). Here we record its key properties.

Given two observable interfaces $P, Q \in \mathcal{M}$, the distinction
functional is
\begin{equation}
  \mathcal{D}(P,Q)
  = H\!\left(\tfrac{P+Q}{2}\right)
    - \tfrac{1}{2}\bigl(H(P) + H(Q)\bigr)
  = \operatorname{JSD}(P\|Q).
\end{equation}
The interface distance is $d_{\mathrm{int}}(P,Q) = \mathcal{D}^{1/2}(P,Q)$.

It is interface-native: its definition requires no reference to any
internal structure. It is symmetric, satisfies the triangle inequality,
and vanishes if and only if $P = Q$. It is the unique metric on
$\mathcal{M}$ that is interface-native, symmetric, and locally equivalent
to the Fisher information metric.

\section{The Derived Interface Conjecture}

The Derived Interface Conjecture is the central unifying claim of the
physical applications of the framework:

\begin{conjecture}[Derived Interface]
Observable spacetime theories are effective interfaces generated by
projection from a larger admissibility manifold. Gauge symmetry, Hilbert
structure, and metric positivity are properties of the interface rather
than of the underlying admissibility structure. Two theories are
physically equivalent if and only if $d_{\mathrm{int}} = 0$ across all
observable energy scales.
\end{conjecture}

In particular, Hilbert-space quantum gravity and Krein-space quadratic
gravity are conjectured to satisfy $d_{\mathrm{int}} = 0$ below the
Planck scale.

%======================================================================
\part{Admissibility}
%======================================================================

\chapter{The Admissibility Program}

\section{Core Idea}

The Admissibility Program asks: which histories, configurations, or
trajectories of a system are physically, cognitively, or institutionally
\emph{realizable}? The answer is not given by a list of allowed states
but by a set of constraint closure conditions: a configuration is
admissible if it satisfies a closed system of mutual constraints.

\begin{definition}[Admissibility Manifold]
An \emph{admissibility manifold} is a triple $(M, \mathcal{A}, \Pi)$
where $M$ is a space of histories, $\mathcal{A} \subseteq M$ is the
admissible subset (those satisfying all constraint closure conditions),
and $\Pi : M \to O$ is the projection to the observable interface.
\end{definition}

The admissible subset $\mathcal{A}$ is determined by the constraint
structure of the theory: which histories satisfy all the equations of
motion, boundary conditions, conservation laws, and causal ordering
conditions simultaneously.

\section{Admissibility Geometry}

The admissibility manifold has a natural geometry: the metric on $M$
that measures how far a history is from the admissibility boundary
$\partial\mathcal{A}$.

\begin{definition}[Admissibility Distance]
The \emph{admissibility distance} of a history $h \in M$ is
\begin{equation}
  d_{\mathcal{A}}(h) = \inf_{h' \in \mathcal{A}} d(h, h'),
\end{equation}
where $d$ is a metric on $M$. A history is admissible iff
$d_{\mathcal{A}}(h) = 0$.
\end{definition}

\begin{definition}[Admissibility Distortion]
The \emph{admissibility distortion} of a theory perturbation $\delta T$
is the change in the admissibility manifold $\mathcal{A}$ induced by
$\delta T$: how the boundary of admissibility shifts.
\end{definition}

\section{Constraint Closure}

A system is \emph{constraint-closed} if its admissibility conditions
are mutually consistent: satisfying one constraint does not force
violation of another. Constraint closure is the condition that the
admissibility manifold $\mathcal{A}$ is non-empty.

Constraint closure failure is the signature of a contradictory or
over-determined theory: a theory that requires more of its histories
than any history can provide. The history of physics contains several
examples: the original formulation of Dirac's equation required both
positive-energy and negative-energy solutions (constraint closure was
restored by the hole theory and later by quantum field theory); the
original formulation of the Einstein--Hilbert action was
non-renormalizable (constraint closure at the quantum level required
additional terms).

\section{Admissibility Equivalence and Theory Space}

Two admissibility manifolds $(M^{(1)}, \mathcal{A}^{(1)}, \Pi_1)$ and
$(M^{(2)}, \mathcal{A}^{(2)}, \Pi_2)$ are \emph{admissibility-equivalent}
if $\Pi_1(\mathcal{A}^{(1)}) = \Pi_2(\mathcal{A}^{(2)})$: their
projections to the observable interface agree on all admissible histories.

This is a weaker condition than structural equivalence (isomorphism of
the internal manifolds) and a stronger condition than empirical equivalence
(agreement on all observed outcomes). Admissibility equivalence is the
correct notion of theoretical equivalence for the purposes of the present
program.

\chapter{Reachability and Continuation}

\section{Reachability Volumes}

The \emph{reachability volume} of a history $h \in \mathcal{A}$ at time
$t$ is the volume of the set of admissible future histories consistent
with $h$ up to time $t$:
\begin{equation}
  \mathrm{Vol}_t(h)
  = \mu\!\bigl(\{h' \in \mathcal{A} : h'|_{[0,t]} = h|_{[0,t]}\}\bigr).
\end{equation}
A history with large reachability volume has many admissible continuations;
a history with small reachability volume is nearly determined by its past.

Reachability volume is the admissibility-theoretic analogue of entropy:
it measures how many futures are open from a given past. The second law
of thermodynamics, in this framework, becomes the claim that reachability
volume is non-decreasing under the macroscopic dynamics, which follows
from the non-injectivity of the coarse-graining projection.

\section{Boundary Effects and the Xylomorphic Criterion}

The boundary of the admissibility manifold $\partial\mathcal{A}$ plays
a special role: histories near the boundary have few admissible
continuations. A system near $\partial\mathcal{A}$ is fragile: small
perturbations may push it outside the admissible region entirely.

The \emph{xylomorphic criterion} is the condition $\lambda < 1$ for
regenerative systems: a system is regenerative (capable of self-repair
and continuation) if its deformation rate is less than its regeneration
rate. Systems satisfying the xylomorphic criterion maintain a positive
distance from $\partial\mathcal{A}$; systems violating it drift toward
the boundary and eventual collapse.

\section{Path Dependence}

The admissibility manifold may be path-dependent: whether a history $h$
is admissible may depend not only on its current state but on how it
arrived there. Path dependence introduces irreversibility into the
admissibility structure itself, not just into the dynamics.

Path dependence is the formal correlate of hysteresis in physics,
institutional lock-in in social theory, and traumatic memory in
cognitive science. In each case, the past constrains the admissible
futures in ways that cannot be fully captured by the current state alone.

\section{Continuation versus Termination}

A \emph{continuation} is a history that extends an existing admissible
history while remaining admissible. A \emph{termination} is a history
that has no admissible continuation: it has reached a boundary of
$\mathcal{A}$ from which there is no return.

The question of whether a given history has continuations is the
admissibility-theoretic version of the survival question: can this
system, process, or institution continue? The theory of continuation
connects the admissibility program to the biology of persistence, the
thermodynamics of dissipative structures, and the political theory of
state legitimacy.

\chapter{Preference Fields}

\section{Preference Geometry}

The Preference Field Program formalizes the geometry of value, preference,
and decision. Its central claim is that preferences are not a list of
desired outcomes but a field defined over the space of possible histories:
a function that assigns a preference intensity to each direction in history
space.

\begin{definition}[Preference Field]
A \emph{preference field} is a smooth function $V : \mathcal{A} \to \mathbb{R}$
on the admissibility manifold, together with a gradient flow
$\nabla V : \mathcal{A} \to T\mathcal{A}$ pointing in the direction of
increasing preference.
\end{definition}

Preference fields generalize utility functions by allowing preferences to
depend on the history of the system rather than only on its current state.
This makes the framework applicable to systems with path-dependent values:
institutions whose goals evolve, cognitive systems whose preferences are
formed by experience, and physical systems whose energy landscape shifts
with their history.

\section{Value Landscapes and Decision Surfaces}

The level sets of the preference field $V$ are \emph{indifference surfaces}:
sets of histories that the system values equally. The gradient of $V$
defines a \emph{decision surface}: the direction in history space that
the system will preferentially explore.

A \emph{local optimum} in the preference field is a history $h^* \in
\mathcal{A}$ such that $\nabla V(h^*) = 0$ and $\nabla^2 V(h^*)$ is
negative definite. It is a history from which any deviation decreases
preference. The distinction between local and global optima is central
to the theory of learning and adaptation: a system that follows local
gradient ascent may become trapped at a local optimum far from the
globally preferred history.

\section{Preference Transport}

As the admissibility manifold changes --- through learning, experience,
or environmental change --- the preference field must be transported
across the changing landscape. This is the problem of \emph{preference
transport}: how do preferences defined relative to one admissibility
structure translate to a new one?

Preference transport is the formal correlate of value alignment in
artificial intelligence, value change in moral philosophy, and preference
revision in decision theory. The tools of parallel transport on manifolds
provide a rigorous framework for this problem: the preference field is
transported along a path in the space of admissibility manifolds, and
the result is a preference field on the new manifold that is as close as
possible to the original in the sense of minimizing the total transported
difference.

%======================================================================
\part{Repair Theory}
%======================================================================

\chapter{Repair as a Primitive}

\section{The Priority of Repair}

Repair is prior to optimization. A system that cannot repair distinctions
that have been collapsed cannot learn from experience, cannot correct
errors, and cannot adapt to novelty. Optimization --- the improvement of
performance along a fixed dimension --- presupposes that the system's
basic distinctions are intact. Repair is the process that keeps them
intact.

Repair is prior to prediction. A prediction requires a stable distinction
between what was predicted and what occurred. When that distinction
collapses --- when the system can no longer tell the difference between
its predictions and its observations --- prediction becomes impossible.
Repair restores the distinction.

Repair is prior to intelligence. Intelligence, in the most general sense,
is the capacity to maintain and extend appropriate distinctions in the
face of environmental pressure. A system that can only optimize cannot
be intelligent in this sense; it can only execute a fixed strategy. A
system that can repair can respond to the unexpected.

\section{Repair versus Optimization}

The distinction between repair and optimization is not a matter of
degree but of kind. Optimization improves performance along a fixed
objective function within a fixed admissibility structure. Repair
restores an admissibility structure that has been damaged or a
distinction space that has been partially collapsed.

\begin{principle}[Repair-Optimization Distinction]
Optimization operates within a fixed admissibility manifold $\mathcal{A}$
and moves toward the maximum of a preference field $V$. Repair operates
on the admissibility manifold itself: it restores $\mathcal{A}$ when it
has been damaged, or restores the distinction space $(X,\sim)$ when
distinctions have been collapsed. Repair changes the structure within
which optimization operates.
\end{principle}

\section{Formal Definition of Repair}

\begin{definition}[Repair Map]
Given a projection setting $(M, O, \Pi)$ and a distinguished admissibility
predicate $\mathcal{A} \subseteq M$, a \emph{repair map} for a history
$h \in M \setminus \mathcal{A}$ is a map $\rho : M \setminus \mathcal{A}
\to \mathcal{A}$ that returns the closest admissible history:
\begin{equation}
  \rho(h) = \operatorname{argmin}_{h' \in \mathcal{A}} d(h, h').
\end{equation}
The \emph{repair cost} of $h$ is $d_{\mathcal{A}}(h) = d(h, \rho(h))$.
\end{definition}

Repair is not in general unique: there may be multiple closest admissible
histories. When repair is non-unique, the system faces a choice about
how to restore admissibility, and that choice reflects the system's
values and preferences.

\section{Repair and the Ecology of Distinctions}

The ecology of a distinction space is the network of relationships
between distinctions: how the preservation or collapse of one distinction
affects others. The \emph{Ecology of Distinctions} (a 31-chapter
monograph) develops this framework in detail.

Key concepts include: distinction \emph{dependencies} (when maintaining
distinction $A$ requires maintaining distinction $B$); distinction
\emph{conflicts} (when maintaining $A$ requires collapsing $C$);
distinction \emph{clusters} (groups of mutually reinforcing distinctions);
and distinction \emph{cascades} (when the collapse of one distinction
triggers the collapse of many others).

The ecology of distinctions provides the framework for understanding
how complex systems --- brains, institutions, languages, ecosystems ---
maintain their coherence against the constant pressure of distinction
collapse.

\chapter{Memory as Repair}

\section{Memory as Distinction Preservation}

Memory is the technology by which a system preserves distinctions through
time. A memory is a record that allows the system to mark what was as
different from what is, and to use that marking in its current operation.

\begin{definition}[Memory]
In the projection framework, a \emph{memory} is a state in the internal
space $M$ that encodes information about past observations: a history
$h_{\text{mem}}$ that is observable-equivalent to the past state
$h_{\text{past}}$ but temporally distinct from it. The \emph{memory
fidelity} is $d(h_{\text{past}}, \rho(h_{\text{mem}}))$, where $\rho$
is the repair map from memory to past state.
\end{definition}

Memory loss is the growth of memory infidelity over time: the increasing
reconstruction defect as the memory trace diverges from the original
state it encodes.

\section{The MEM|8 Framework}

MEM|8 (Wave-Stabilized Memory Architecture) is a formal framework for
persistent memory in systems with dynamic substrates. Its core insight
is that memory is stabilized not by the permanence of the substrate but
by wave interference: memories are encoded in stable interference patterns
in a dynamic medium, analogous to holograms.

Key components of MEM|8:
\begin{enumerate}[label=(\roman*)]
  \item \emph{Event logs}: explicit records of distinction-preserving
        transitions in the history of the system.
  \item \emph{Ecphory}: the process of retrieving a memory by providing
        partial information (a cue) and completing the distinction pattern.
  \item \emph{Persistent histories}: the representation of memory as a
        history in the internal space $M$ rather than as a state.
  \item \emph{Memory reconstruction}: the inference of the generating
        history from partial memory traces, using repair maps to fill
        in collapsed distinctions.
\end{enumerate}

MEM|8 is authored by C.\ and A.\ Chenoweth.

\section{Ecphory and Distinction Completion}

Ecphory is the memory process of retrieving a full distinction pattern
from a partial cue. In distinction geometry terms, it is a repair
operation: the cue provides a partial specification of the target
distinction pattern, and ecphory completes the pattern to the nearest
admissible full specification.

The quality of ecphory depends on the structure of the distinction space:
if the target pattern is in a dense region of $\mathcal{A}$ (many nearby
admissible patterns), ecphory will find it readily; if it is in a sparse
region, the repair map will produce a nearby but not identical result.

\chapter{Intelligence as Distinction Repair}

\section{The Distinction-Repair Account of Intelligence}

Intelligence, in the most general sense, is the capacity to maintain and
extend appropriate distinctions in the face of environmental pressure
that would collapse them. This account subsumes:

\textit{Understanding} as the capacity to maintain the distinctions
relevant to a domain: to keep separate what should be kept separate and
to identify as the same what should be identified as the same.

\textit{Learning} as the capacity to update the distinction space in
response to experience: to repair collapsed distinctions, to refine
coarse distinctions, and to create new distinctions where none existed.

\textit{Scientific discovery} as the distinction reorganization that makes
previously invisible differences visible: the introduction of a framework
in which things that seemed the same are distinguished, and things that
seemed different are unified.

\textit{Error correction} as the repair of collapsed distinctions in a
formal system: the restoration of the distinction between correct and
incorrect, true and false, valid and invalid.

\textit{Anomaly integration} as the process of incorporating observations
that violate current distinctions into a revised distinction space: the
repair of the distinction between observation and expectation.

\section{Learning as Projection Revision}

Learning is, in the CLIO framework, the revision of the projection $\pi :
M_{\mathrm{int}} \to M_{\mathrm{obs}}$. A learning event occurs when the
current projection fails --- when the current internal model predicts an
observation that does not occur, or fails to predict one that does.
Learning is the update of $\pi$ (or of the internal model $M_{\mathrm{int}}$)
to reduce this mismatch.

Gradient descent is one mechanism for this update: adjust $\pi$ in the
direction that reduces the discrepancy between $\pi(m)$ and the observed
$o$. But gradient descent is not the only mechanism. Paradigm shifts,
creative leaps, and abductive inference are also learning mechanisms, and
they typically involve a more radical revision of the distinction space
than gradient descent can achieve: they are distinction reorganizations
rather than distinction refinements.

%======================================================================
\part{Coordination and Society}
%======================================================================

\chapter{Coordination Geometry: The HYDRA Framework}

\section{Overview}

The HYDRA framework (Coordination Geometry and Collective Admissibility)
formalizes the geometry of coordination: the conditions under which
multiple agents with potentially distinct internal models and distinction
spaces can act coherently.

The central insight of HYDRA is that coordination is not a property of
individual agents but of the \emph{colimit} of their admissibility
manifolds: the minimal admissibility structure that contains all
individual structures and is consistent with their interactions.

\begin{definition}[Collective Admissibility]
Given a collection of agents $(M_i, \mathcal{A}_i, \Pi_i)_{i \in I}$
and interaction maps $\phi_{ij} : M_i \to M_j$, the \emph{collective
admissibility manifold} is the categorical colimit
\begin{equation}
  \mathcal{A}_{\mathrm{coll}}
  = \operatorname{colim}_{i \in I} \mathcal{A}_i,
\end{equation}
computed in the category of measurable spaces with the interaction maps
as morphisms. An action profile $(h_i)_{i \in I}$ is \emph{collectively
admissible} if the induced element in $\mathcal{A}_{\mathrm{coll}}$ is
admissible.
\end{definition}

\section{Synchronization and the Kuramoto Model}

The Kuramoto model of coupled oscillators provides the simplest example
of collective admissibility geometry. Each oscillator $i$ has a phase
$\theta_i \in [0, 2\pi)$ and a natural frequency $\omega_i$. The
interaction among oscillators is given by the coupling
\begin{equation}
  \dot\theta_i = \omega_i + \frac{K}{N}\sum_{j=1}^N \sin(\theta_j - \theta_i),
\end{equation}
where $K$ is the coupling strength. The collectively admissible states
are the synchronized states in which all oscillators have the same
effective frequency.

In HYDRA terms, each oscillator has an admissibility manifold
$\mathcal{A}_i = [0,2\pi)$ and the collective admissibility manifold
is the subspace of phase profiles consistent with synchronization. The
transition from incoherence to synchronization at the critical coupling
$K_c$ is a phase transition in the collective admissibility geometry.

\section{Social Coherence and Institutional Alignment}

The HYDRA framework extends the Kuramoto intuition to institutions.
An institution is a collection of agents with individual admissibility
manifolds (individual roles, constraints, and objectives) that interact
through interaction maps (reporting structures, information flows,
incentive systems).

Institutional coherence is the analog of synchronization: the condition
in which the individual admissibility manifolds are aligned enough that
collectively admissible action profiles exist. Institutional dysfunction
is incoherence: the collective admissibility manifold is empty or very
sparse, and coordinated action becomes impossible.

\chapter{Institutions as Distinction Maintenance}

\section{Institutions as Distinction-Preservation Technologies}

An institution is, at its most basic, a technology for maintaining
distinctions that would otherwise collapse. Legal systems maintain the
distinction between permitted and forbidden. Market systems maintain the
distinction between prices and quantities. Scientific institutions
maintain the distinction between evidence and conjecture. Democratic
institutions maintain the distinction between legitimate and illegitimate
power.

The health of an institution is therefore measurable as its capacity
to maintain the distinctions it was designed to preserve. Institutional
decay is distinction collapse. Institutional corruption is distinction
confusion: the inability to maintain the distinction between the
institution's purpose and the interests of those who operate it.

\section{Legibility and the Hilbert Assumption in Institutions}

Scott's notion of \emph{legibility} --- the process by which states
simplify and standardize the societies they govern --- is the institutional
form of the Hilbert Assumption: the identification of the institution's
representation of a community with the community itself.

When a state maps a complex, evolving, locally structured community onto
a legible grid --- standard plots, surnames, tax categories, occupational
classifications --- it creates a simplified model whose projection
onto observables (tax records, census data, administrative outcomes) is
tractable. But the simplification collapses distinctions that matter to
the people in the community, and the resulting policies, calibrated to the
simplified model rather than to the community, produce characteristic
pathologies.

The framework identifies this as a projection failure: the state's model
$\pi(m)$ mismatches the community's observable behavior $o$, not because
the community is irrational but because the state's projection $\pi$ does
not track the relevant distinctions.

\section{Goodhart Effects as Hidden Admissibility Directions}

Goodhart's Law --- when a measure becomes a target, it ceases to be a
good measure --- is, in the framework, an instance of the hidden
admissibility direction phenomenon. When an institution optimizes a
proxy measure $M$ for a true objective $T$, it explores a fiber of
the projection $\pi : \text{behavior space} \to \text{measure space}$
that keeps $M$ fixed while changing $T$. If the fiber is non-trivial ---
if behaviors exist that achieve high $M$ while degrading $T$ ---
optimization along $M$ will find those behaviors.

The remedy is to tighten the projection: to redesign the measurement
system so that the fiber of $\pi^{-1}(\text{high } M)$ contains only
behaviors that also achieve high $T$. This is the institutional analog
of removing hidden admissibility directions: making the relevant behavior
distinctions visible in the measurement system.

\chapter{Civilization as a Repair System}

\section{Large-Scale Distinction Preservation}

Civilization, at the largest scale, is a system for preserving
distinctions against the constant pressure of collapse. The distinctions
civilizations preserve include: the distinction between written and lost
knowledge (libraries, education systems, scholarly traditions); the
distinction between living and dead languages; the distinction between
enforced and unenforced law; the distinction between calibrated and
uncalibrated measurement; and the distinction between functional and
dysfunctional infrastructure.

Each of these is maintained by an active technology: not just a record
but a practice of maintaining the record, using it, updating it, and
repairing it when it is damaged. Civilization is not an archive. It is
a repair system for the distinctions that make further repair possible.

\section{Science as Distinction Accumulation}

Science accumulates distinctions: it creates and preserves the distinctions
that allow the world to be described with increasing precision. The
distinction between mass and weight; between heat and temperature; between
correlation and causation; between statistical significance and practical
significance --- each of these is a hard-won distinction, created by
scientific work and maintained by scientific practice.

Scientific institutions are distinction-maintenance technologies: journals,
peer review, replication, and systematic criticism are all mechanisms for
preventing the collapse of scientific distinctions into confusion, bias,
or fashion.

%======================================================================
\part{Computation}
%======================================================================

\chapter{Computation as Restriction of Possibility}

\section{The Reachability Account of Computation}

Computation, in the reachability framework, is the process of restricting
the space of admissible futures. A computation begins with a large
admissibility manifold (many possible output states) and proceeds by
applying constraint after constraint, each one eliminating inadmissible
futures, until only the desired output remains.

\begin{principle}[Computational Restriction]
Computation is not the creation of outputs but the elimination of
inadmissible alternatives. A correct computation is one that restricts
the admissibility manifold to exactly the correct answer.
\end{principle}

This account is the computational analog of the physical insight that
constraint-first dynamics is more fundamental than force-first dynamics.
In physics, the laws of motion are not forces that push states forward
but constraints that select admissible trajectories from all possible
ones. In computation, algorithms are not procedures that create outputs
but constraint systems that select the admissible output from all possible
symbol strings.

\section{Irreversible Operations and Distinction Collapse}

Irreversible logical operations --- AND, OR, NAND --- collapse distinctions:
they map multiple input states to a single output state. The AND gate maps
$(0,0)$, $(0,1)$, and $(1,0)$ all to $0$, collapsing the distinction
between them. The energy cost of this collapse is Landauer's limit:
$k_BT\ln 2$ per bit erased.

Reversible computation preserves distinctions: every input state maps to
a distinct output state, so the computation is injective. Reversible
computation has no Landauer cost, but it requires more hardware: the
intermediate states that irreversible computation discards must be stored
in reversible computation.

The distinction geometry framework gives a unified account of both:
irreversible computation increases reconstruction defect (the output
does not determine the input), while reversible computation preserves
it. The tradeoff between energy cost and information retention is the
tradeoff between accepting reconstruction defect and paying the energy
cost to avoid it.

\chapter{Flow Computing: The NCL Framework}

\section{Null Convention Logic}

Null Convention Logic (NCL), developed by Karl Fant, is a formalism for
asynchronous, self-timed digital circuits. Its key property is that
computation is \emph{complete}: the circuit indicates when it has
produced a valid output and when it is resetting, eliminating the need
for a global clock.

In distinction geometry terms, NCL is a computational formalism in which
the circuit maintains the distinction between ``computing'' and
``complete'' at all times. The absence of a global clock means that the
circuit's operation is not synchronized to an external distinction
(the clock cycle) but to an internal one (the completion signal).

NCL circuits are history-first: the output of a gate is not determined
by the current input alone but by the history of inputs, in particular
by whether the circuit has received a complete, valid input pattern since
its last reset. This makes NCL circuits natural candidates for
computation on dynamic substrates where the timing of inputs is not
controlled.

\section{Dynamic Substrates}

A \emph{dynamic substrate} is a computational medium whose structure
changes during computation: its admissibility manifold shifts as the
computation proceeds. Neural tissue is a dynamic substrate: synaptic
weights change in response to activity, altering the admissibility
conditions for future computation. An economy is a dynamic substrate:
the rules of exchange, the availability of resources, and the legal
framework all change in response to economic activity.

Computation on dynamic substrates requires a formalism in which the
distinction space $(X, \sim)$ is not fixed but evolves. The CLIO
projection framework handles this by allowing the projection $\pi$ to
change: $\pi_t : M_{\mathrm{int}} \to M_{\mathrm{obs}}$ is a
time-dependent projection whose evolution encodes the changing
admissibility structure.

\chapter{Programs as Histories: The Spherepop Framework}

\section{Overview}

Spherepop is an irreversible event calculus in which the fundamental
objects are not states but histories. A Spherepop program is a record
of events --- Pop, Refuse, Bind, Collapse --- each of which modifies
the admissibility manifold of the computation.

\begin{definition}[Spherepop Operators]
\begin{enumerate}[label=(\roman*)]
  \item \textbf{Pop}: eliminates the distinction between a suspended
        event and its continuation, resolving a pending computation.
  \item \textbf{Refuse}: maintains the distinction between the current
        state and an inadmissible continuation, blocking a computation
        that would violate admissibility.
  \item \textbf{Bind}: creates a distinction between two events by
        linking them: one can only occur if the other has occurred.
  \item \textbf{Collapse}: irreversibly eliminates all distinctions
        within a set of alternatives, selecting one and discarding the
        rest.
\end{enumerate}
\end{definition}

\section{Histories versus States}

The distinction between history-based and state-based computation is
central to the Spherepop framework. In state-based computation, the
computation is described by a function from states to states: the output
is determined by the current state alone. In history-based computation,
the output is determined by the entire history of the computation.

History-based computation is more expressive: it can represent
dependencies that are invisible in state-based formulations. In
particular, Refuse --- the operator that blocks inadmissible continuations
--- cannot be expressed in a purely state-based formalism, because the
admissibility of a continuation may depend on the history of how the
current state was reached.

Spherepop is implemented as a C interpreter; the formal semantics is
developed in the framework of irreversible event calculi.

%======================================================================
\part{Dynamics and Substrates}
%======================================================================

\chapter{Frozen Processes}

\section{When Moving Variables Become Hidden Constants}

A \emph{frozen process} is a dynamic process in which one or more
variables that should be treated as evolving have been fixed at a
particular value. The fixing may be deliberate (a model simplification)
or inadvertent (a failure to notice that the variable changes).

The frozen-process critique is a recurring theme in the Flyxion program.
It arises whenever a model replaces a trajectory with a state: whenever
what is actually a history is described as a snapshot, and the dynamics
of the history is attributed to the snapshot itself.

\begin{example}[Frozen Time in Physics]
The Hamiltonian formalism of classical mechanics ``freezes'' time: it
describes the evolution of a system by a vector field on phase space,
treating time as a parameter rather than a dynamical variable. This
works well for conservative systems but breaks down when the system's
constraints are time-dependent. The frozen-time approximation hides the
time-dependence in the constraints, making it invisible to methods that
only look at the phase space.
\end{example}

\begin{example}[Frozen Preferences in Economics]
Standard economic models freeze preferences: they treat the utility
function as a fixed parameter of the agent. But preferences change in
response to experience, advertising, habituation, and social influence.
Freezing preferences hides the dynamics of preference formation, making
it impossible to model how markets shape the values of their participants.
\end{example}

\section{Static Models and Their Failure Modes}

Static models --- models that describe a system by its state rather than
its history --- fail when the dynamics of the system's history matters
for its current behavior. The failure modes are characteristic:

\textit{Path dependence} appears as irreducible complexity: two systems
with the same current state behave differently because they arrived there
by different paths. Static models cannot represent this.

\textit{Hysteresis} appears as inexplicable history-dependence: a system
that behaves differently after undergoing a cycle than it did before.
Static models can only represent hysteresis by introducing additional
state variables, effectively recovering a history-based description.

\textit{Emergence} appears as behavior that cannot be predicted from the
current state of the components: it requires knowing how the components
interacted over time. Static models cannot represent true emergence.

\chapter{Dynamic Substrates}

\section{Computing on Evolving Manifolds}

When the admissibility manifold $\mathcal{A}$ changes over time, both
the valid computations and the valid results of those computations change.
Computing on a dynamic substrate requires tracking the evolution of
$\mathcal{A}$ as well as the state of the computation.

The TARTAN framework (Trajectory-Aware Recursive Tiling with Annotated
Noise) addresses this problem. TARTAN represents the admissibility manifold
as a tiling of the history space, with each tile annotated by the
admissibility conditions that hold within it. As the manifold evolves,
the tiling is updated, and the annotations reflect the new conditions.

\section{Plastic Systems}

A system is \emph{plastic} if its admissibility manifold changes in
response to its own history: the constraints that govern the system's
behavior evolve as the system acts. Neural tissue is plastic: synaptic
plasticity changes the admissibility conditions for neural computation
in response to activity. Economic systems are plastic: the legal and
institutional framework that governs economic activity changes in response
to economic outcomes.

Plasticity is the computational correlate of learning: a plastic system
updates its own constraint structure in response to experience, thereby
changing what is admissible for it in the future. The theory of plasticity
is therefore the computational version of the theory of learning, and the
distinction geometry of plastic systems is the formal framework within
which both can be understood.

\chapter{Distinction Reorganization}

\section{Conceptual Change and Paradigm Shifts}

A \emph{distinction reorganization} is a change in the distinction space
$(X, \sim)$ that is neither a refinement nor a coarsening: it introduces
new distinctions and collapses old ones simultaneously. Scientific
paradigm shifts are the most dramatic examples, but distinction
reorganizations occur at every scale of cognitive and social life.

\begin{example}[The Copernican Reorganization]
The Copernican revolution replaced a distinction space in which the
distinction between sublunary and superlunary was fundamental with one
in which the distinction between planets and fixed stars replaced it.
The old distinction (earthly versus heavenly) was collapsed; the new
one (planetary versus stellar motion) was created. This is a genuine
distinction reorganization, not a refinement or coarsening.
\end{example}

\section{Creativity as Distinction Creation}

Creativity is, in the framework, the production of new distinctions where
none previously existed: the identification of a difference that had not
previously been marked, or the introduction of a concept that allows
things to be distinguished that were previously conflated.

The geometry of creative acts is the geometry of paths in the space of
distinction spaces: a creative act is a movement to a point in the space
of equivalence relations on $X$ that was previously unoccupied. The
difficulty of creativity is that the space of distinction spaces is vast,
and most movements in it are arbitrary. Creative acts are movements to
points in this space that are both novel and useful: that create
distinctions with high downstream value for repair, coordination,
and prediction.

%======================================================================
\part{Physical Applications}
%======================================================================

\chapter{RSVP: The Relativistic Scalar-Vector Plenum}

\section{Overview}

RSVP (Relativistic Scalar-Vector Plenum) is a field-theoretic cosmology
grounded in constraint-first dynamics. Its central claim is that
cosmological structure --- the distribution of matter, the geometry of
spacetime, the arrow of time --- emerges from the constraint dynamics of
three coupled fields: a scalar field $\Phi$, a vector field $\mathbf{v}$,
and an entropy density $S$.

The RSVP framework is a physical instantiation of the distinction geometry
program: the three fields encode, respectively, the intensity, transport,
and entropy of distinction-preserving structure in the cosmological medium.

\section{The Field Equations}

The RSVP field equations govern the coupled dynamics of $(\Phi, \mathbf{v}, S)$:
\begin{align}
  \partial_\mu\partial^\mu\Phi + \xi R\Phi &= J[\mathbf{v}, S], \\
  \nabla \cdot \mathbf{v} + \partial_t S &= \Sigma[\Phi, \mathbf{v}], \\
  \partial_t S + \mathbf{v}\cdot\nabla S &= \Pi[\Phi, \mathbf{v}].
\end{align}
Here $R$ is the Ricci scalar, $J$, $\Sigma$, and $\Pi$ are coupling
functionals, and $\xi$ is the conformal coupling constant. The scalar
$\Phi$ drives the geometry through its stress-energy tensor; the vector
$\mathbf{v}$ transports distinction-preserving structure; the entropy
$S$ measures the local degree of distinction collapse.

\section{The RSVP Transport Operator and the Born Rule}

The RSVP transport operator is
\begin{equation}
  Z(x) = \Phi(x)\, e^{\widehat{L}(x)},
  \qquad
  \widehat{L}(x) = \mathbf{v}(x)\cdot\nabla + \theta(x)\widehat{T},
\end{equation}
where $\Phi$ supplies scaling, $\mathbf{v}\cdot\nabla$ supplies transport,
and $\theta\widehat{T}$ supplies torsional structure. The observable
transition probability between configurations $C_i$ and $C_j$ is
\begin{equation}
  T_{ij} = \int_{\gamma : C_i \to C_j} |Z[\gamma]|^2\,\mathcal{D}\gamma.
\end{equation}
When the coarse-graining projection has fibers admitting a unitary action
preserved by RSVP transport, the induced transition matrix is unistochastic:
$T_{ij} = |U_{ij}|^2$ for a unitary $U$. The Born rule of quantum
mechanics is recovered as the leading-order approximation when the RSVP
dynamics reduces to linear field theory.

\section{Cosmological Implications}

In the cosmological context, RSVP provides a framework for deriving
large-scale structure from constraint dynamics rather than from initial
conditions. The observed homogeneity, isotropy, and flatness of the
universe emerge as consequences of admissibility-flow consistency: the
configurations that maximize reachability volume are those of maximum
symmetry.

The entropy field $S$ provides a direct connection to the arrow of time:
the direction in which $S$ increases is the direction in which
distinction-preserving structure is being consumed by the dynamics.
This makes the arrow of time an emergent property of the RSVP constraint
structure rather than a fundamental asymmetry of the laws of nature.

\chapter{Cosmology}

\section{Constraint Geometry in Cosmology}

The large-scale geometry of the universe --- its homogeneity, isotropy,
and spatial flatness --- is a constraint geometry problem: which
cosmological histories satisfy all the observational and physical
constraints simultaneously?

The gravitational entropy program of Boyle and Turok approaches this
question from within the projection framework: the observed universe is
a typical element of the admissibility manifold of cosmological histories,
where typicality is measured by the volume of the admissibility set
consistent with the observations. Cosmological simplicity is not a
fine-tuned initial condition but a consequence of the admissibility
geometry.

\section{The CPT-Symmetric Universe}

The CPT-Symmetric Universe proposal (Boyle, Finn, Turok) provides a
projection-based account of the Big Bang. The universe is CPT-symmetric:
the post-bang universe is the CPT image of the pre-bang universe, and
both are determined by a CPT-invariant vacuum state at the singularity.

In projection terms, the CPT symmetry is a constraint on the admissibility
manifold of cosmological histories: only those histories that are CPT-symmetric
are admissible. This constraint selects a unique vacuum state, eliminates
the need for fine-tuned initial conditions, and determines the density
perturbation spectrum from the CPT-symmetric quantum vacuum.

\section{Entropy and Expansion}

The expansion of the universe is, in the RSVP framework, the growth of
reachability volume: as the universe expands, the number of admissible
future histories grows. The cosmological arrow of time is the direction
of reachability growth.

This connects RSVP cosmology to the distinction geometry account of
irreversibility: the expansion of the universe is a distinction-creation
process, not a distinction-collapse process. The universe creates new
distinctions (new regions of spacetime, new physical structures) as it
expands; the local thermodynamic arrow of time, which is a
distinction-collapse process, is a consequence of local dynamics within
an expanding background that is creating distinctions globally.

\chapter{Quantum Interfaces}

\section{Projection Geometry in Quantum Mechanics}

The quantum mechanical measurement process is a projection from the
internal state space (the Hilbert space $\mathcal{H}$ or Krein space
$\mathcal{K}$) to the observable space (the spectrum of measurement
outcomes). The Born rule is the formula for computing the probability
distribution on the observable space induced by this projection.

The central claim of the quantum interface program is that the Born rule
is not a fundamental axiom but a consequence of the projection structure.
The squaring operation $a \mapsto |a|^2$ is a projection from complex
amplitude space to non-negative intensity space; the phase is discarded
because it is a hidden direction in the fiber of this projection.

\section{The Ghost Problem as a Projection Question}

The ghost problem in higher-derivative quantum gravity is, in the
projection framework, the question of whether negative-norm states in
the internal state space $M_{\mathrm{adm}}$ are hidden admissibility
directions or observable pathologies.

A negative-norm state is a hidden admissibility direction if the
projection $\Pi : M_{\mathrm{adm}} \to M_{\mathrm{obs}}$ maps its
contribution to zero: if the observable probability distribution is
unaffected by the presence of the ghost. The Turok--Bateman proposal
--- formulated in the Krein-space framework with ghost-parity symmetry
--- is a claim that the massive ghost graviton of quadratic gravity is a
hidden admissibility direction of exactly this kind.

\section{Krein Spaces and the Observable Interface}

A Krein space $\mathcal{K} = \mathcal{K}_+ \oplus \mathcal{K}_-$ with
indefinite inner product is a natural $M_{\mathrm{adm}}$ for a quantum
theory with ghost states. The observable interface is the pushforward
distribution on the space of observable outcomes:
\begin{equation}
  P(F \mid I)
  = \mathrm{tr}_+\!\bigl((\Pi_F S \Pi_I)^\dagger (\Pi_F S \Pi_I)\bigr),
\end{equation}
where the trace is taken with respect to the positive Hilbert space
structure on $\mathcal{K}_+$. When ghost-parity symmetry holds
($[G,S]=0$), this probability functional is positive and normalised,
even though the underlying space $\mathcal{K}$ has indefinite inner
product.

\chapter{Derived Interface Physics}

\section{Observable Physics as Projection}

The derived interface program is the application of the projection
hierarchy to the full structure of physics. Its central claim is that
the structural features of physical theories --- Hilbert space, gauge
symmetry, metric positivity, the Born rule, and spacetime geometry itself
--- are properties of the observable interface rather than of the
underlying admissibility manifold.

\begin{conjecture}[Derived Interface Physics]
Observable spacetime physics is the effective interface generated by
projecting an underlying admissibility manifold $M_{\mathrm{adm}}$
onto the observable domain $M_{\mathrm{obs}}$. The hierarchy of effective
theories --- quantum mechanics, quantum field theory, general relativity,
quadratic gravity --- represents successive approximations to the
projection $\Pi$ at successively higher orders in the admissibility-flow
expansion.
\end{conjecture}

\section{Gauge Symmetry as Fiber Structure}

Gauge symmetry is, in the projection framework, fiber structure: the
gauge orbit through a field configuration is the fiber $\Pi^{-1}(\Pi(A))$,
the set of all configurations that project to the same observable
outcome. The gauge symmetry of the theory is the symmetry of the fiber.

This gives gauge symmetry a projective interpretation rather than a
fundamental one: gauge symmetry is not a property of the physical world
but of the representational choice that introduces more degrees of freedom
than the observables require. The physical content of the theory is
the projection, not the gauge structure.

\section{The Admissibility-Flow Expansion}

Spacetime curvature, in the derived interface program, is the first-order
correction to admissibility flow: a measure of how much the projection
$\Pi$ departs from a globally flat map. The Einstein field equations are
the conditions for admissibility-flow consistency at leading order.
Quadratic gravity is the second-order term in the expansion. Higher-derivative
theories represent higher-order terms, with ghost sectors corresponding
to negative-metric directions in the fiber of $\Pi$ that do not reach
the observable interface.

%======================================================================
\part{Meta-Theory}
%======================================================================

\chapter{Against Latent Fundamentalism}

\section{The Central Error}

Latent fundamentalism is the error of granting ontological status to
internal coordinates merely because they appear in a successful model.
It is the confusion of the instrument with the reality it instruments:
the map with the territory, the model with the world, the projection
machinery with the observable world it projects onto.

The error is natural and persistent. Every successful model creates a
temptation: if this model works, perhaps the objects it refers to are
the real objects of the world. The wave function works, so perhaps the
wave function is real. The utility function works, so perhaps utility is
real. The organizational chart works, so perhaps the organization really
is organized that way.

The Flyxion program resists this temptation not by denying that internal
structures exist but by insisting on the distinction between what is
observable and what is internal. The internal structure is a tool for
generating and explaining observable distinctions. It is indispensable.
But it is not ontologically primary.

\section{A Catalogue of Latent Fundamentalist Errors}

The following is a partial catalogue of latent fundamentalist errors,
each analyzed in terms of the projection framework:

\textit{Wave function realism}: treating the wave function as a physical
object rather than as an element of the internal space $M_{\mathrm{adm}}$
from which observable probabilities are projected.

\textit{Utility realism}: treating the utility function as a real feature
of agents rather than as a projection of their revealed preferences onto
a mathematical representation.

\textit{Organizational chart realism}: treating the organizational chart
as the institution rather than as a projection of actual decision-making
and power structures.

\textit{Metric realism in quantum gravity}: treating the positive-definiteness
of the Hilbert inner product as a necessary feature of any physically
admissible theory rather than as a sufficient condition for observable
consistency.

\textit{GDP fundamentalism}: treating GDP as a measure of welfare rather
than as a projection of economic activity onto a single number that
collapses the distinction between types of activity that matter for welfare
in very different ways.

\section{The Correction}

The correction for latent fundamentalism is always the same: identify the
projection $\Pi$, locate the observable interface $M_{\mathrm{obs}}$, and
restrict admissibility conditions to that interface. Internal coordinates
are then free to be whatever they need to be to generate the correct
observable predictions, without those coordinates being awarded
ontological status.

This is not instrumentalism: the claim is not that internal structure is
unreal or that we should be agnostic about it. The claim is that internal
structure is epistemologically secondary: the appropriate object of inquiry
is the geometry of the observable interface, and internal structures are
evaluated by their utility for generating and explaining that geometry.

\chapter{Distinction Geometry}

\section{The Interface-Native Metric}

Distinction Geometry is the formal development of the geometry of
distinguishability between observable interfaces. Its primitive is the
Jensen--Shannon distance
\begin{equation}
  d_{\mathrm{int}}(P,Q)
  = \operatorname{JSD}^{1/2}(P\|Q)
  = \left[H\!\left(\tfrac{P+Q}{2}\right)
    - \tfrac{1}{2}\bigl(H(P)+H(Q)\bigr)\right]^{1/2},
\end{equation}
which compares observable interfaces without reference to any underlying
admissibility manifold.

The distinction functional $\mathcal{D}(P,Q) = \operatorname{JSD}(P\|Q)$
measures the information destroyed when the distinction between $P$ and
$Q$ is forgotten: the entropy of the merged interface minus the average
entropy of the separate interfaces. This is the information that
disappears when an observer stops distinguishing between the two theories.

\section{Theory Space and Its Curvature}

The space of theories $(\mathcal{M}, d_{\mathrm{int}})$ is a metric
space that is locally Riemannian with the Fisher information metric.
Its curvature has a precise interpretation in terms of the difficulty
of theoretical inference:

Positive sectional curvature at a theory $P$ means that nearby theories
are more similar in their observable predictions than their parameter-space
distance suggests: many distinct internal models collapse onto nearly
identical observables. This is the geometry of overfitting, paradigm
stability, and institutional opacity.

Negative sectional curvature at a theory $P$ means that small changes
in the theory produce large changes in observable predictions: the theory
is highly falsifiable in that parameter direction.

Zero curvature is the Euclidean regime: parameter distance and observable
distance scale proportionally, and the theory is locally well-identified.

\section{Reconstruction Defect and Irreversibility}

The reconstruction defect $\Delta_R(h) = d(h, R(\Pi(h)))$ measures the
minimum error that any reconstruction must make in recovering the generating
history from an observation. When the projection $\Pi$ is non-injective,
no reconstruction map achieves zero defect: irreversibility is a theorem
about the structure of projection, not a contingent feature of dynamics.

This unifies the arrow of time, thermodynamic irreversibility, memory loss,
scientific underdetermination, and institutional goal drift as instances
of a single phenomenon: persistent reconstruction defect under non-injective
projection.

\chapter{The Ecology of Distinctions}

\section{Overview}

The Ecology of Distinctions is a 31-chapter academic monograph (current
length: approximately 30,000 lines of \LaTeX) that develops the full
formal theory of how distinctions interact, support each other, compete
with each other, and form systems. It is the most detailed formal
development of the distinction program.

The monograph covers: formal distinction theory and its connection to
equivalence relations; reachability in distinction spaces; admissibility
geometry; the boundary theory of distinction collapse; repair algebras
for the restoration of collapsed distinctions; distinction ecology
(the network structure of distinction dependencies and conflicts);
biological applications (cell types as distinction spaces, morphogenesis
as distinction creation, homeostasis as distinction preservation); and
a full treatment of the admissibility manifold hierarchy including proofs
of the Projection Pathology Separation Theorem and the Observable
Admissibility Theorem.

\section{The Master Chain}

The Ecology of Distinctions concludes with the following chain, which
is the organizing principle of the entire program:
\begin{equation}
  \text{Distinction}
  \;\to\;
  \text{Information}
  \;\to\;
  \text{Entropy}
  \;\to\;
  \text{Repair}
  \;\to\;
  \text{Continuation}
  \;\to\;
  \text{Admissibility}
  \;\to\;
  \text{Worlds.}
\end{equation}
Distinctions enable information. Information enables entropy accounting.
Entropy accounting enables repair. Repair enables continuation. Continuation
generates admissibility conditions. Admissibility conditions, accumulated
over histories, generate the structured world.

%======================================================================
\part{Unified Framework}
%======================================================================

\chapter{The Master Dependency Chain}

\section{From Distinctions to Worlds}

The architectures described in the preceding parts are not independent
inventions. Each is a level in a single dependency chain. The chain
goes:

\medskip
\noindent
\textbf{Distinctions} are the primitive. A system that can tell things
apart can do everything else that follows. A system that cannot is
indistinguishable, by any observation, from a system with no internal
structure at all.

\medskip
\noindent
\textbf{Information} is maintained distinction. The Shannon entropy of
a source is the expected distinguishability it produces. Mutual information
is the distinctions shared between two sources. Compression is
distinguishability reduction. Information theory is the theory of
distinctions in the presence of uncertainty.

\medskip
\noindent
\textbf{Entropy} is collapsed distinction. Thermodynamic entropy measures
the number of microstates compatible with a macrostate --- the size of
the fiber of the coarse-graining projection. High entropy means many
internal histories produce the same observable, which means many
distinctions have been collapsed by the projection. The second law is
the statement that, under the dynamics, the fiber grows.

\medskip
\noindent
\textbf{Repair} is distinction restoration. When a distinction has been
collapsed --- when the fiber has grown, when information has been lost,
when the gap between the internal history and its observable image has
increased --- repair is the attempt to restore the collapsed distinction
from whatever information remains. Repair is possible only to the extent
that the collapsed distinction left traces in the observable world; it
is impossible to the extent that the collapse was complete.

\medskip
\noindent
\textbf{Continuation} is the capacity of a system to extend its history
into the future while remaining admissible. A system that has lost too
many distinctions --- whose reachability volume has collapsed too far ---
cannot continue: it has reached a boundary of $\mathcal{A}$ from which
there is no admissible exit. Repair maintains the distinctions that make
continuation possible.

\medskip
\noindent
\textbf{Admissibility} is the accumulated constraint closure of a
history: the set of conditions that a history must satisfy to count as
realizable. Admissibility is generated by the interaction of distinctions
over time: each distinction imposes constraints on which subsequent
distinctions are possible. The admissibility manifold is the set of all
histories that satisfy all constraints simultaneously.

\medskip
\noindent
\textbf{Worlds} are admissibility manifolds that have achieved the kind
of closure, continuity, and coherence that we associate with a persistent,
interacting domain of reality. A world, in this sense, is a
distinction-preservation system that has become self-sustaining: its
distinctions generate the constraints that maintain those distinctions.

\section{The Role of Each Framework}

Within this chain, each framework in the Flyxion program occupies a
specific role:

\textbf{RSVP} instantiates the chain in a physical field theory: $\Phi$
carries distinction intensity, $\mathbf{v}$ carries distinction transport,
and $S$ carries distinction entropy. The RSVP field equations are the
physical laws of distinction flow.

\textbf{CLIO} formalizes the projection layer: it is the theory of how
internal distinction spaces map onto observable ones, and what can be
inferred about the former from the latter.

\textbf{TARTAN} handles computation on dynamic substrates: the case where
the admissibility manifold changes during the process of distinction
maintenance.

\textbf{HYDRA} extends the framework to collectives: the case where
multiple distinction spaces must be coordinated into a collective
admissibility manifold.

\textbf{MEM|8} formalizes the memory layer: the technology by which
a system preserves distinctions through time against the pressure of
collapse.

\textbf{Spherepop} provides the computational semantics: a formalism
for processes that are irreversible, history-dependent, and
distinction-collapsing.

\textbf{The Admissibility Program} develops the geometry of admissibility
manifolds: the formal theory of which histories are realizable and how
the space of realizable histories is structured.

\textbf{Repair Theory} develops the algebra of distinction restoration:
the formal theory of how collapsed distinctions are partially restored
from surviving traces.

\textbf{Distinction Geometry} provides the interface metric: the formal
theory of how observable interfaces are compared, how theories are
distinguished from each other, and what the geometry of that comparison
looks like.

\textbf{Coordination Geometry} extends the projection framework to
multi-agent systems: the formal theory of when agents with different
distinction spaces can achieve collective admissibility.

\textbf{Preference Fields} provide the value layer: the formal theory
of how preferences over distinctions are represented, transported, and
traded off.

\chapter{Open Problems}

\section{Mathematical Open Problems}

The following mathematical problems arise directly from the frameworks
documented in this work and have not yet been fully resolved:

\subsection*{Distinction Geometry}

\textbf{Completeness of the theory space.} Is the metric space
$(\mathcal{M}, d_{\mathrm{int}})$ complete? The question is whether
Cauchy sequences of observable interfaces converge to an interface that
can be realised by some admissibility manifold.

\textbf{Uniqueness of the JSD metric.} Is $d_{\mathrm{int}} = \operatorname{JSD}^{1/2}$
the unique interface-native metric satisfying symmetry, triangle inequality,
and local Fisher-metric equivalence? An axiomatic characterization would
clarify its status as the canonical metric on theory space.

\textbf{Curvature bounds.} What bounds on the sectional curvature of
$(\mathcal{M}, g^F)$ are implied by admissibility conditions? In
particular, does high admissibility (a small admissibility manifold
relative to the full history space) imply high curvature of the theory
space?

\textbf{Repair algebras.} What algebraic structures do repair maps
form? In particular: is the composition of two repair maps a repair
map? Is there an identity repair? What is the appropriate notion of
morphism between repair systems?

\subsection*{RSVP and Physics}

\textbf{Unistochastic emergence.} Under what conditions does RSVP
transport generate a unistochastic transition kernel? The current
conjecture (fibers admitting a $\mathrm{U}(n)$ action preserved by the
dynamics) needs to be verified for the specific transport operators
that arise in RSVP cosmology.

\textbf{Ghost-parity in the spin-2 sector.} The most pressing open
problem in the physical applications is whether the ghost-parity symmetry
of the Turok--Bateman Krein-space formulation of quadratic gravity extends
to the spin-2 sector. If it does, the Derived Interface Conjecture gains
strong support. If it does not, the conjecture must be revised.

\textbf{RSVP field equations and general relativity.} Under what
conditions do the RSVP field equations reduce to the Einstein field
equations in an appropriate limit? The leading-order admissibility-flow
expansion is conjectured to yield general relativity, but the precise
conditions need to be established.

\subsection*{Coordination and Collective Dynamics}

\textbf{HYDRA colimit existence.} Under what conditions does the
categorical colimit $\mathcal{A}_{\mathrm{coll}} = \operatorname{colim}_i
\mathcal{A}_i$ exist and is non-empty? This is the formal question of
when coordination is possible at all.

\textbf{Phase transitions in collective admissibility.} The transition
from incoherence to synchronization in the Kuramoto model is a phase
transition in the collective admissibility geometry. What are the
analogous phase transitions in institutional systems, and can they be
predicted from the geometry of the individual admissibility manifolds
and the interaction maps?

\section{Physical Open Problems}

\textbf{RSVP and the cosmological constant.} The gravitational entropy
program (Boyle and Turok) predicts that the observed value of the
cosmological constant is typical for a universe with the observed
entropy budget. Can this prediction be made quantitative from within
the RSVP framework?

\textbf{Dark matter and the CPT universe.} The CPT-symmetric universe
predicts that the stable right-handed neutrino has mass
$M \approx 5 \times 10^8\,\mathrm{GeV}$ to account for the observed
dark matter density. This prediction will be testable by future
large-scale structure surveys.

\textbf{Primordial gravitational waves.} The CPT-symmetric universe
predicts no primordial tensor modes ($r = 0$ at leading order). Future
experiments (CMB-S4, LiteBIRD) with sensitivity $r \lesssim 0.003$
will test this prediction.

\textbf{Asymptotic freedom and the Higgs as a composite.} The
dimension-zero scalar program (Boyle and Turok) predicts that the Higgs
is composite, with mass set by the dynamical scale of the asymptotically
free conformal sector. The mechanism is analogous to dimensional
transmutation in QCD. What is the predicted Higgs compositeness scale?

\section{Conceptual Open Problems}

\textbf{The boundary between repair and learning.} When does a system
repair a collapsed distinction (restoration of what was) versus learn a
new distinction (creation of something new)? The formal distinction
between repair and learning in the distinction geometry framework has
not been made precise.

\textbf{Distinction creation and the origin of novelty.} The framework
accounts for the preservation, collapse, repair, and reorganization of
distinctions. It does not yet give a full account of the creation of
genuinely new distinctions: distinctions for which no prior distinction
served as a template. This is the formal version of the problem of
genuine novelty.

\textbf{The geometry of institutional failure.} The framework identifies
institutional opacity, goal drift, and Goodhart effects as instances of
hidden admissibility directions and distinction collapse. But it does
not yet give a quantitative theory of institutional failure: at what
rate do institutions drift, and what determines the rate?

\textbf{Consciousness as distinction maintenance.} If consciousness is
the capacity to maintain distinctions --- to be aware of what is different
from what --- then the framework provides the beginning of a formal
account of consciousness. But the relationship between distinction
maintenance and phenomenal experience is not yet worked out. This is
the hardest problem in the program.

\newpage

\section*{Appendices}
\appendix

\chapter{Notation Reference}

\begin{tabular}{p{3.5cm}p{8cm}}
\toprule
Symbol & Meaning \\
\midrule
$(X, \sim)$ & Distinction space: set $X$ with equivalence relation $\sim$ \\
$[x]_\sim$ & Equivalence class of $x$ \\
$X/{\sim}$ & Quotient space of equivalence classes \\
$\Pi : M \to O$ & Projection from history space to observable space \\
$F_o = \Pi^{-1}(o)$ & Fiber over observable $o$ \\
$\mathcal{A} \subseteq M$ & Admissibility set \\
$d_\mathcal{A}(h)$ & Admissibility distance of history $h$ \\
$\Delta_R(h)$ & Reconstruction defect of $h$ under map $R$ \\
$\delta(\Pi)$ & Minimum worst-case reconstruction defect \\
$\mathcal{D}(P,Q)$ & Distinction functional $= \operatorname{JSD}(P\|Q)$ \\
$d_{\mathrm{int}}(P,Q)$ & Interface distance $= \operatorname{JSD}^{1/2}(P\|Q)$ \\
$(\mathcal{M}, d_{\mathrm{int}})$ & Theory space \\
$g^F_{ij}$ & Fisher information metric \\
$H(P)$ & Shannon entropy of distribution $P$ \\
$M_{PQ}$ & Equal-weight mixture $\tfrac{1}{2}(P+Q)$ \\
$Z(x)$ & RSVP transport operator $\Phi(x)e^{\widehat{L}(x)}$ \\
$\widehat{L}$ & Transport generator $\mathbf{v}\cdot\nabla + \theta\widehat{T}$ \\
$T_{ij}$ & Observable transition matrix \\
$\mathcal{K} = \mathcal{K}_+ \oplus \mathcal{K}_-$ & Krein space decomposition \\
$G = \Pi_+ - \Pi_-$ & Ghost parity operator \\
$[u,v]_\mathcal{K}$ & Krein inner product \\
\bottomrule
\end{tabular}

\chapter{Index of Frameworks}

\begin{description}[style=nextline]
  \item[Admissibility Program]
    Distinguishability geometry and reachability-based ontology.
    Chapters 7--8, Parts~III and~VIII.

  \item[CLIO (Constraint-Layered Inference and Projection)]
    Formal theory of projection from internal to observable models.
    Chapter~5, Part~II.

  \item[Coordination Geometry / HYDRA]
    Collective admissibility as categorical colimit; multi-agent coordination.
    Chapter~13, Part~V.

  \item[Distinction Geometry]
    Jensen--Shannon interface metrics and the curvature of distinguishability.
    Chapter~27, Part~IX.

  \item[Ecology of Distinctions]
    Network theory of distinction dependencies, conflicts, and cascades.
    Chapter~28, Part~IX; also a 31-chapter standalone monograph.

  \item[MEM|8 (Wave-Stabilized Memory Architecture)]
    Formal framework for persistent memory in dynamic substrates.
    Chapter~11, Part~IV. Authored by C.\ and A.\ Chenoweth.

  \item[Preference Fields]
    Geometry of value, preference, and decision over admissibility manifolds.
    Chapter~9, Part~III.

  \item[Repair Theory]
    Formal theory of distinction restoration; repair as a cognitive primitive.
    Chapters~10--12, Part~IV.

  \item[RSVP (Relativistic Scalar-Vector Plenum)]
    Field-theoretic cosmology grounded in constraint-first dynamics.
    Chapter~22, Part~VIII.

  \item[Spherepop]
    Irreversible event calculus with Pop, Refuse, Bind, and Collapse operators.
    Chapter~18, Part~VI. Implemented as a C interpreter.

  \item[TARTAN (Trajectory-Aware Recursive Tiling with Annotated Noise)]
    Computation on dynamic substrates with evolving admissibility manifolds.
    Chapter~20, Part~VII.
\end{description}

\newpage

%----------------------------------------------------------------------
\begin{thebibliography}{99}

\bibitem{flyxion_ecology}
Flyxion,
\textit{The Ecology of Distinctions},
31-chapter monograph on distinction theory, reachability, and admissibility
geometry (in preparation).

\bibitem{flyxion_cpr}
Flyxion,
\textit{Constraint, Projection, and Reachability} (CPR),
92-chapter LuaLaTeX monograph (in preparation).

\bibitem{flyxion_quadratic}
Flyxion,
``Quadratic Gravity as a Derived Interface Theory:
Admissibility Manifolds, Projection Geometry, and the Emergence of
Observable Physics,''
preprint (2025).

\bibitem{flyxion_distinction_geometry}
Flyxion,
\textit{Distinction Geometry: Projection, Reconstruction Defect, and the
Interface-Native Metric} (Mathematical Core),
preprint (2025).

\bibitem{boyle2018cpt}
L.~Boyle, K.~Finn, and N.~Turok,
``CPT-symmetric universe,''
\textit{Physical Review Letters} \textbf{121}, 251301 (2018).

\bibitem{turokboyle2022entropy}
N.~Turok and L.~Boyle,
``Gravitational entropy and the flatness, homogeneity and isotropy puzzles,''
\textit{Physics Letters B} \textbf{849}, 138442 (2024).

\bibitem{stelle1977}
K.~S. Stelle,
``Renormalization of higher-derivative quantum gravity,''
\textit{Physical Review D} \textbf{16}, 953--969 (1977).

\bibitem{endres2003}
D.~M. Endres and J.~E. Schindelin,
``A new metric for probability distributions,''
\textit{IEEE Transactions on Information Theory} \textbf{49}(7),
1858--1860 (2003).

\bibitem{amari2000}
S.-I. Amari and H.~Nagaoka,
\textit{Methods of Information Geometry},
American Mathematical Society, Providence, 2000.

\bibitem{needham1997}
T.~Needham,
\textit{Visual Complex Analysis},
Oxford University Press, Oxford, 1997.

\bibitem{barandes2023}
J.~A. Barandes,
``The unistochastic quantum theory,''
\textit{arXiv:2302.10778} [quant-ph] (2023).

\bibitem{scott1998}
J.~C. Scott,
\textit{Seeing Like a State},
Yale University Press, New Haven, 1998.

\bibitem{kuhn1962}
T.~S. Kuhn,
\textit{The Structure of Scientific Revolutions},
University of Chicago Press, Chicago, 1962.

\bibitem{landauer1961}
R.~Landauer,
``Irreversibility and heat generation in the computing process,''
\textit{IBM Journal of Research and Development} \textbf{5}(3),
183--191 (1961).

\bibitem{jaynes2003}
E.~T. Jaynes,
\textit{Probability Theory: The Logic of Science},
Cambridge University Press, Cambridge, 2003.

\bibitem{fant1996}
K.~M. Fant,
\textit{Logically Determined Design: Clockless System Design with NULL
Convention Logic},
Wiley-Interscience, Hoboken, 2005.

\bibitem{kuramoto1975}
Y.~Kuramoto,
``Self-entrainment of a population of coupled non-linear oscillators,''
in \textit{International Symposium on Mathematical Problems in Theoretical
Physics}, Lecture Notes in Physics~39, Springer, Berlin, 1975, pp.~420--422.

\bibitem{kay2023}
J.~Kay and M.~King,
\textit{Radical Uncertainty},
Bridge Street Press, London, 2020.

\end{thebibliography}

\end{document}
